\section{Discussione}
\label{sec:discussione}

\subsection{Che cosa dicono i risultati}
\label{sec:discussione-risultati}

Il confronto fra solver e programmi zero-dimensionali mostra che il modello è
stato inserito correttamente nel solver: le due strade, che fanno la stessa
fisica, danno le stesse curve entro lo 0{,}2\% in quasi tutti i casi, anche
con le reazioni accese. Il confronto con il lavoro di riferimento mostra che
anche la fisica è quella giusta. Nei casi senza reazioni le temperature finali
coincidono con quelle imposte dalla conservazione dell'energia e con i valori
del lavoro di riferimento entro lo 0{,}3\%, cioè entro l'incertezza con cui si
leggono le figure; nei casi con reazioni le temperature a fine calcolo
coincidono entro lo 0{,}8\%, la composizione della miscela di azoto entro il
2{,}5\% e quella dell'aria entro il 15\%. Nel transitorio l'accordo
è di pochi punti percentuali per l'azoto puro, per la miscela di azoto e
ossigeno e per l'aria. La differenza più grande riguarda la temperatura
vibrazionale nelle miscele di azoto molecolare e atomico, dove arriva al
10--20\%.

\subsection{L'origine delle differenze}
\label{sec:discussione-origine}

L'origine più probabile del ritardo della temperatura vibrazionale nelle
miscele è il modo in cui \mpp{} media i tempi di rilassamento relativi ai
diversi partner di collisione, che non è quello del lavoro di riferimento. Il
modo di mediare conta solo quando i partner sono più di uno: nell'azoto puro
l'accordo è buono, mentre nelle miscele con azoto atomico la temperatura
vibrazionale sale più tardi. Lo stato finale non ne risente, perché dipende
solo dall'energia. L'ipotesi non è stata verificata implementando la media del
lavoro di riferimento, che resta uno degli sviluppi possibili. Per l'azoto puro
e per la miscela di azoto e ossigeno restano piccoli spostamenti nel tempo, di
qualche punto percentuale, la cui origine non è stata isolata.

\subsection{Che cosa insegna il confronto}
\label{sec:discussione-lezioni}

Buona parte del lavoro è consistita nel ricostruire scelte che l'articolo non
dichiara, o dichiara in modo incompleto: se l'energia elettronica è accesa,
quale energia guida lo scambio vibrazionale, con quale esponente si calcola la
temperatura delle dissociazioni, quali costanti di velocità sono usate per
l'aria. Nessuna di queste scelte è un dettaglio: sbagliarne una sposta i
risultati molto più delle differenze rimaste, come mostrano i confronti della
sezione~\ref{sec:risultati}, con scarti fino al 56\% sulla temperatura
vibrazionale, un fattore 3 sulle densità dell'aria e circa il 25\% sulla
temperatura finale. Lo stesso articolo ne dà un esempio quando spiega la
differenza con LeMANS, nel caso della miscela di azoto molecolare e atomico,
con una diversa convenzione nella
correzione di Park.

Per prendere queste decisioni due strumenti sono stati decisivi. Il primo è
la temperatura finale imposta dalla conservazione dell'energia, che non
dipende da nessun modello di rilassamento e dice subito se i dati del gas sono
quelli giusti. Il secondo sono le curve estratte dalle figure vettoriali, che
trasformano il confronto con l'articolo da una lettura a occhio in una misura.
Infine, la verifica su due strade separa gli errori di implementazione dalle
scelte di modello: poiché solver e programmi coincidono, ogni differenza
rimasta rispetto al lavoro di riferimento viene dal modello, non da un errore
nel codice.
